\chapter{Die Ligaseite, von oben nach unten}\label{kap:ligaseite}

Eine Ligaseite (\code{render\_league\_page},
\datei{RCode/generate_static_site.R:516--641}) besteht aus diesen Teilen:

\begin{enumerate}
  \item Kopf mit Wortmarke und Navigation (Kapitel~\ref{kap:seiten}),
  \item Hinweis auf veraltete Daten (verborgen, solange die Daten frisch sind),
  \item \textbf{Prognose}: Heatmap und zwei Zusammenfassungstabellen,
  \item \textbf{Tabelle}: Ligatabelle mit ELO, bei Regionalligen mit Zonen,
        darunter der Link zum ELO-Verlauf,
  \item \textbf{Rückblick}: die Spiele des aktuellen Spieltags mit Prognose und
        ELO-Änderung,
  \item \textbf{Live}: laufende Spiele, nur wenn welche laufen,
  \item \textbf{Ausblick}: die kommenden Spiele mit Wahrscheinlichkeiten und
        Ergebnis-Matrix,
  \item Fuß mit Zeitstempel und Blog-Link.
\end{enumerate}

\screenshot{ligaseite-bundesliga}{Kopf, Heatmap und Zusammenfassungstabellen der Bundesliga-Seite.}

\section{Prognose: die Heatmap}

Die Sektion beginnt mit \enquote{Saisonprognose \emph{Liga}} und dem Satz
\enquote{10\,000 Monte-Carlo-Simulationen des Saisonrests. Jede Zeile zeigt,
mit welcher Wahrscheinlichkeit ein Team auf welchem Platz landet.} Die
Heatmap ist die Matrix $\hat P$ aus Gleichung~\eqref{eq:mc}. Jede Zeile ist
eine Mannschaft, sortiert nach erwartetem Platz, jede Spalte ein Endplatz.
Die Zeilenköpfe zeigen das Kürzel. Der volle Vereinsname erscheint als
Tooltip, beim Antippen auch auf Touch-Geräten.

\subsection*{Die Farbformel}

Jede Zelle wird von Weiß zum Tintenblau \code{\#1F3A4D} interpoliert
(\code{.heat\_style}, \datei{RCode/render_helpers.R:83--96}). Damit kleine
Wahrscheinlichkeiten sichtbar bleiben und nicht linear verblassen, wird die
Wahrscheinlichkeit vorher mit dem Exponenten 0{,}6 gestaucht:
\begin{equation}
  t = p^{\,0{,}6}, \qquad
  \begin{pmatrix} R\\ G\\ B \end{pmatrix}
  = \operatorname{round}\!\left(
    \begin{pmatrix} 255\\ 255\\ 255 \end{pmatrix}
    + t\left[\begin{pmatrix} 31\\ 58\\ 77 \end{pmatrix} - \begin{pmatrix} 255\\ 255\\ 255 \end{pmatrix}\right]
  \right).
\end{equation}
Ist $t > 0{,}52$, also ab $p \approx 33{,}6\,\%$, wechselt die Schrift von
fast Schwarz (\code{\#15130F}) zu Weiß.

\begin{center}
\definecolor{h01}{RGB}{241,243,244}
\definecolor{h10}{RGB}{199,206,210}
\definecolor{h50}{RGB}{107,125,138}
\definecolor{h100}{RGB}{31,58,77}
\begin{tabular}{@{}lcccc@{}}
\toprule
$p$ & 1\,\% & 10\,\% & 50\,\% & 100\,\% \\
$t = p^{0{,}6}$ & 0{,}063 & 0{,}251 & 0{,}660 & 1 \\
RGB & (241,\,243,\,244) & (199,\,206,\,210) & (107,\,125,\,138) & (31,\,58,\,77) \\
Farbe &
\colorbox{h01}{\makebox[2.6em]{\ttfamily\small 1}} &
\colorbox{h10}{\makebox[2.6em]{\ttfamily\small 10}} &
\colorbox{h50}{\makebox[2.6em]{\ttfamily\small\color{white} 50}} &
\colorbox{h100}{\makebox[2.6em]{\color{white}\checkmark}} \\
\bottomrule
\end{tabular}
\end{center}

\subsection*{Was in einer Zelle steht}

Eine Zelle mit $p = 0$ bleibt leer und weiß. Sonst steht dort der
Prozentwert nach der Regel aus Abschnitt~\ref{sec:prozent}, ohne
Prozentzeichen, denn die Legende sagt \enquote{Angaben in Prozent}. Auf
schmalen Bildschirmen (bis 520~px) werden Werte unter einem Prozent nur noch
als Färbung gezeigt. Die Heatmap passt so ohne seitliches Scrollen auf ein
Telefon.

\section{Prognose: die Zusammenfassungstabellen}

Unter der Heatmap stehen zwei Tabellen nebeneinander, oben die Spitze und
unten das Tabellenende. Jede Spalte summiert eine Gruppe von Plätzen.
Gezeigt werden nur Mannschaften, deren Summe über die Spalten mindestens
1~Prozent beträgt. Eine 0 erscheint hier als \enquote{0} und nicht als
leere Zelle wie in der Heatmap.

\begin{table}[ht]
\centering\small
\caption{Spaltengruppen je Liga (\texttt{RCode/league\_views.R}). Negative
Zahlen zählen von unten. \enquote{berechnet} heißt, dass der Wert nicht aus
einer Platzsumme stammt, sondern aus Kapitel~\ref{kap:regionalliga}.}
\label{tab:panels}
\begin{tabularx}{\textwidth}{@{}l>{\raggedright}X>{\raggedright\arraybackslash}X@{}}
\toprule
Liga & oben & unten \\
\midrule
Bundesliga & Meister (1), Champions League (2--4), Europa League (5), Conference League Quali (6) & Relegation (16), Abstieg (17--18) \\
2.~Bundesliga & Aufstieg (1--2), Relegation Bundesliga (3) & Relegation 3.~Liga (16), Abstieg (17--18) \\
3.~Liga & aus dem Malus-Lauf: Aufstieg (1--2), Relegation (3), DFB-Pokal (4) & Abstieg (17--20) \\
Frauen-Bundesliga & Meister (1), Champions League Quali (2--3) & Abstieg ($-2$ bis $-1$) \\
2.~Frauen-Bundesliga & aus dem Malus-Lauf: Aufstieg (1--2) & Abstieg ($-3$ bis $-1$) \\
RL Nord, RL Bayern & Meister (1), Aufstieg (berechnet, Aufstiegsspiele) & Abstieg (berechnet); Bayern zusätzlich Relegation \\
RL Nordost, West, SüdWest & Aufstieg (1, Direktaufstieg) & Abstieg (berechnet) \\
\bottomrule
\end{tabularx}
\end{table}

Berechnete Spalten werden über ihren Namen zugeordnet, nicht über ihre
Position. Fehlt eine, bricht das Rendern dieser Seite ab, statt eine falsche
Zahl zu zeigen (\datei{generate_static_site.R:246--285}).

\section{Tabelle: Ligatabelle und ELO}

Die Sektion \enquote{Ligatabelle und ELO} zeigt den aktuellen Stand, daneben
die Stärkeschätzung des Modells und ihre Veränderung seit Saisonbeginn
(\code{render\_liga\_tabelle}, \datei{RCode/render_sections.R:115--164}).

\begin{center}
\small
\begin{tabular}{@{}lll@{}}
\toprule
Spalte & Inhalt & sortierbar \\
\midrule
Platz    & Tabellenplatz & ja \\
Team     & voller Name, Kürzel als Tooltip & -- \\
Sp.      & Spiele (auf Telefonen ausgeblendet) & -- \\
Tordiff. & Tordifferenz mit Vorzeichen (auf Telefonen ausgeblendet) & -- \\
Punkte   & Punkte & ja \\
ELO      & aktuelle ELO, eine Nachkommastelle & ja \\
$\Delta$ ELO & ELO heute minus ELO zum Saisonstart & ja \\
\bottomrule
\end{tabular}
\end{center}

Die Tabelle zählt beendete und gewertete Spiele und sortiert nach Punkten,
Tordifferenz und Toren. Ein Klick auf einen Spaltenkopf sortiert im Browser
um. Die Sortierung nach $\Delta$ ELO zeigt oben, wer die Erwartung vom
Saisonstart am deutlichsten übertrifft (Issue~\#285). Vorzeichen werden vom
\emph{gerundeten} Wert bestimmt: \enquote{$+$}, \enquote{$-$} (echtes
Minuszeichen U+2212) oder \enquote{$\pm$} bei null.

\subsection*{Auf- und Abstiegszonen der Regionalligen}

In den Regionalligen markiert eine dünne farbige Linie am linken Zeilenrand
die Zonen: \textcolor{akzent}{Rot} für Abstieg, \textcolor{ocker}{Ocker} für
Relegation (nur Bayern), \textcolor{gruen}{Grün} für Aufstieg. Die
\emph{Deckkraft} der Linie ist die Wahrscheinlichkeit selbst, mindestens aber
15~Prozent, damit jede Chance über null sichtbar bleibt. Treffen mehrere
Zonen zusammen, gilt der Vorrang Rot vor Ocker vor Grün. Die Zonen gehören
zum Tabellenplatz und verschwinden, sobald nach einer anderen Spalte sortiert
wird.

Darunter erklärt eine Fußnote die Regel der Staffel, etwa für Nord:
\enquote{Drei Vereine steigen ab, je Absteiger aus der 3.~Liga einer mehr.
Gewinnt Nord das Aufstiegsspiel, ist es einer weniger.} Schwankt die Zahl
der Absteiger, folgt ihre Verteilung (\enquote{3 mit $x$\,\%, 4 mit
$y$\,\% …}) und danach die Wahrscheinlichkeit je Platz.

Unter der Tabelle führt ein Link zur ELO-Verlaufsseite der Liga
(Abschnitt~\ref{sec:verlauf}).

\section{Rückblick}

Der Rückblick zeigt die beendeten Spiele des aktuellen Spieltags. Solange
dieser nicht abgeschlossen ist, zeigt er zusätzlich den vorigen und dazu
gekennzeichnete Nachholspiele. Jede Zeile hat fünf Spalten:

\begin{enumerate}
  \item Anstoß in Berliner Zeit, etwa \enquote{Sa. 4.10., 15:30 Uhr},
  \item die Paarung \enquote{Heim – Gast},
  \item einen 1/X/2-Balken mit den Wahrscheinlichkeiten \emph{vor} dem Spiel:
        \textcolor{tinte}{Heimsieg}, \textcolor{ink3}{Remis} (hellgrau),
        \textcolor{ocker}{Auswärtssieg},
  \item das Ergebnis,
  \item die ELO-Änderung als \enquote{Heim / Gast}, etwa
        \enquote{$+4{,}1$ / $-4{,}1$}.
\end{enumerate}

Die Wahrscheinlichkeiten des Balkens stammen aus Abschnitt~\ref{sec:tor-beispiel}
und werden mit der ELO unmittelbar vor dem Spiel berechnet. Gerundet wird so,
dass die drei Segmente zusammen genau 100 ergeben:
\begin{equation}
  h = \operatorname{round}(100\,p_H), \qquad a = \operatorname{round}(100\,p_A), \qquad x = 100 - h - a.
\end{equation}
Segmente unter 8~Prozent tragen keine Zahl. Gewertete Spiele sind mit
\enquote{Wertung} markiert und haben keine ELO-Änderung.

\section{Live und Ausblick}

\textbf{Live} erscheint nur, wenn Spiele laufen. Die Sektion zeigt Paarung
und Zwischenstand mit dem Hinweis \enquote{Prognosen werden während des
Spiels nicht aktualisiert.}

\textbf{Ausblick} zeigt die offenen Spiele des aktuellen und des nächsten
Spieltags, gerechnet mit der ELO von heute, jeweils mit 1/X/2-Balken. Zu
jeder Paarung lässt sich die \emph{Ergebnis-Matrix} aufklappen
(\datei{render_sections.R:561--587}). Sie zeigt das $7\times7$-Gitter aus
Kapitel~\ref{kap:tormodell}: Zeilen sind die Tore der Heimmannschaft, Spalten
die der Gastmannschaft, und die letzte Zeile und Spalte fassen \enquote{6+}
zusammen. Die Werte haben eine Nachkommastelle, Zellen unter 0{,}1~Prozent
bleiben leer. Die Färbung ist relativ zum größten Wert der Matrix:
\code{.heat\_style}$\bigl(\min(1,\ 0{,}75 \cdot p / p_{\max})\bigr)$. Auch
das wahrscheinlichste Ergebnis wird also nie ganz dunkel.

\section{Fuß}

Links steht \enquote{Letztes Update: TT.MM.JJJJ HH:MM MESZ} in Berliner Zeit,
dahinter der maschinenlesbare Zeitstempel für den Veraltet-Hinweis, bei
gedrosseltem Takt zusätzlich \enquote{Eingeschränkter Service}. Rechts steht
\enquote{Mehr dazu unter 30punkte.wordpress.com}.
